\chapter{Taylor Formulas and Asymptotic Expansions}\label{ch:b1:taylor}

Near a point, a smooth function is as good as a polynomial --- with a
controllable error. The Taylor formulas make this exact in three
flavors (integral remainder, Lagrange remainder, Young remainder),
and the resulting \emph{asymptotic expansions}, manipulated
algebraically, become the sharpest tool of elementary analysis:
limits, equivalents, local behavior, asymptotes.

\section{Comparison notation}

\begin{definition}[Landau notation]\label{def:b1:taylor:landau}
Let $f, g$ be defined near $x_0$ ($x_0 \in \R$ or $\pm\infty$). One
writes, as $x \to x_0$:\index{Landau notation}
\begin{itemize}
  \item $f = o(g)$ (``little-o'') when $f = \varepsilon g$ with
        $\varepsilon(x) \to 0$;
  \item $f = O(g)$ (``big-O'') when $f = u g$ with $u$ bounded near
        $x_0$;
  \item $f \sim g$ (``equivalent''\index{equivalent functions}) when
        $f = (1 + \varepsilon) g$ with $\varepsilon \to 0$ ---
        equivalently $f - g = o(g)$.
\end{itemize}
The same notation applies to sequences ($n \to \infty$).
\end{definition}

\begin{proposition}[Rules]\label{prop:b1:taylor:landaurules}
As $x \to x_0$:
\begin{enumerate}
  \item $\sim$ is an equivalence relation; $f \sim g$ implies that
        $f$ and $g$ share limits, signs (near $x_0$), and zeros'
        absence;
  \item equivalents \emph{multiply and divide}: $f_1 \sim g_1$, $f_2
        \sim g_2$ imply $f_1 f_2 \sim g_1 g_2$ and $\frac{f_1}{f_2}
        \sim \frac{g_1}{g_2}$;
  \item equivalents do \emph{not} add: $x + 1 \sim x$ and $-x \sim
        -x + 2$ at $+\infty$, yet the sums $1$ and $2$ are not
        equivalent. To add, return to expansions with explicit
        $o(\cdot)$ terms;
  \item $o(g) + o(g) = o(g)$, $\;u \cdot o(g) = o(ug)$, $\;o(o(g)) =
        o(g)$, and $f \sim g \iff f = g + o(g)$.
\end{enumerate}
\end{proposition}

\begin{proof}
Each is a one-line manipulation of the definitions; for instance
$f_1 f_2 = (1+\varepsilon_1)(1+\varepsilon_2) g_1 g_2$ and $(1 +
\varepsilon_1)(1+\varepsilon_2) \to 1$. The counterexample in (3) is
the proof of (3).
\end{proof}

\section{The three Taylor formulas}

\begin{theorem}[Taylor with integral remainder]\label{thm:b1:taylor:integral}
Let $f$ be of class $C^{n+1}$ on an interval containing $a$ and $x$.
Then\index{Taylor formula}
\[
f(x) = \sum_{k=0}^{n} \frac{f^{(k)}(a)}{k!}\,(x - a)^k
+ \int_a^x \frac{(x - t)^n}{n!}\, f^{(n+1)}(t)\, \dd t .
\]
\end{theorem}

\begin{proof}
Induction on $n$. For $n = 0$: $f(x) = f(a) + \int_a^x f'(t)\dd t$ is
the fundamental theorem (\cref{thm:b1:integration:ftc}). Step:
integrate the remainder by parts,
\[
\int_a^x \frac{(x-t)^n}{n!} f^{(n+1)}(t)\,\dd t
= \Bigl[-\frac{(x-t)^{n+1}}{(n+1)!} f^{(n+1)}(t)\Bigr]_a^x
+ \int_a^x \frac{(x-t)^{n+1}}{(n+1)!} f^{(n+2)}(t)\,\dd t ,
\]
the bracket contributing the term $\frac{f^{(n+1)}(a)}{(n+1)!}(x -
a)^{n+1}$.
\end{proof}

\begin{theorem}[Taylor--Lagrange inequality]\label{thm:b1:taylor:lagrange}
Let $f$ be $C^{n+1}$ with $\abs{f^{(n+1)}} \leq M$ between $a$ and
$x$. Then
\[
\Bigl| f(x) - \sum_{k=0}^{n} \frac{f^{(k)}(a)}{k!}(x-a)^k \Bigr|
\leq M\, \frac{\abs{x - a}^{n+1}}{(n+1)!} .
\]
\end{theorem}

\begin{proof}
Bound the integral remainder: $\bigl|\int_a^x \frac{(x-t)^n}{n!}
f^{(n+1)}(t)\,\dd t\bigr| \leq M \bigl|\int_a^x
\frac{\abs{x-t}^n}{n!}\dd t\bigr| = M\frac{\abs{x-a}^{n+1}}{(n+1)!}$.
\end{proof}

\begin{theorem}[Taylor--Young]\label{thm:b1:taylor:young}
Let $f$ be $n$ times differentiable at $a$. Then, as $x \to a$:
\[
f(x) = \sum_{k=0}^{n} \frac{f^{(k)}(a)}{k!}\,(x-a)^k
+ o\bigl((x-a)^n\bigr) .
\]
\end{theorem}

\begin{proof}
Induction on $n$. For $n = 1$ this is the definition of the
derivative (\cref{def:b1:derivative:def}). Assume the statement at
order $n - 1$, and let $f$ be $n$ times differentiable at $a$. Apply
the induction hypothesis to $f'$ (which is $n-1$ times differentiable
at $a$):
\[
f'(t) = \sum_{k=0}^{n-1} \frac{f^{(k+1)}(a)}{k!}(t-a)^k + r(t),
\qquad r(t) = o\bigl((t-a)^{n-1}\bigr).
\]
Let $g(x) = f(x) - \sum_{k=0}^{n} \frac{f^{(k)}(a)}{k!}(x-a)^k$; then
$g' = r$ and $g(a) = 0$. Given $\varepsilon > 0$, choose $\delta$
with $\abs{r(t)} \leq \varepsilon\abs{t - a}^{n-1}$ for $\abs{t-a}
\leq \delta$; the mean value inequality
(\cref{thm:b1:derivative:mvt}) applied on the segment from $a$ to
$x$ (where $\abs{g'} \leq \varepsilon\abs{x-a}^{n-1}$) yields
$\abs{g(x)} \leq \varepsilon\abs{x - a}^n$: exactly $g(x) =
o((x-a)^n)$.
\end{proof}

\begin{proposition}[Standard expansions at $0$]\label{prop:b1:taylor:standard}
As $x \to 0$, for every fixed order $n$:\index{Taylor
expansion!standard}
\begin{align*}
\eu^x &= 1 + x + \frac{x^2}{2!} + \dots + \frac{x^n}{n!} + o(x^n),\\
\cos x &= 1 - \frac{x^2}{2!} + \frac{x^4}{4!} - \dots
+ \frac{(-1)^p x^{2p}}{(2p)!} + o(x^{2p+1}),\\
\sin x &= x - \frac{x^3}{3!} + \dots + \frac{(-1)^p
x^{2p+1}}{(2p+1)!} + o(x^{2p+2}),\\
\frac{1}{1 - x} &= 1 + x + x^2 + \dots + x^n + o(x^n),\\
\ln(1 + x) &= x - \frac{x^2}{2} + \frac{x^3}{3} - \dots +
\frac{(-1)^{n-1} x^n}{n} + o(x^n),\\
(1 + x)^\alpha &= 1 + \alpha x + \frac{\alpha(\alpha-1)}{2!}x^2 +
\dots + \binom{\alpha}{n} x^n + o(x^n),
\end{align*}
where $\binom{\alpha}{n} = \frac{\alpha(\alpha - 1)\cdots(\alpha - n
+ 1)}{n!}$ for real $\alpha$. ($\cosh$ and $\sinh$: same as $\cos$,
$\sin$ without the alternating signs.)
\end{proposition}

\begin{proof}
Each function is smooth near $0$ with derivatives easy to evaluate:
$(\eu^x)^{(k)} = \eu^x$; the derivatives of $\sin$ and $\cos$ cycle
with period $4$; $\bigl((1+x)^\alpha\bigr)^{(k)} = \alpha(\alpha -
1)\cdots(\alpha - k + 1)(1+x)^{\alpha - k}$;
$\bigl(\ln(1+x)\bigr)^{(k)} = \frac{(-1)^{k-1}(k-1)!}{(1+x)^k}$.
Apply Taylor--Young at $a = 0$. (The geometric one is exact:
$\frac{1}{1-x} - \sum_0^n x^k = \frac{x^{n+1}}{1 - x} = o(x^n)$.)
\end{proof}

\begin{method}[Computing with expansions]\label{met:b1:taylor:compute}
\begin{enumerate}
  \item \emph{Fix the target order} $n$ first, and truncate every
        intermediate result there --- carrying higher terms is
        wasted work, dropping lower ones is an error.
  \item \emph{Sums, products}: expand each factor to order $n$ and
        multiply, discarding beyond $x^n$.
  \item \emph{Composition} $f(u(x))$ with $u(x) \to 0$: substitute
        the expansion of $u$ into that of $f$, order by order.
  \item \emph{Quotients}: write $\frac{1}{1 + v}$ with $v \to 0$ and
        use the geometric expansion.
  \item \emph{Integrate} an expansion term by term (differentiating
        requires more care --- justification: integral of $o(t^n)$
        from $0$ to $x$ is $o(x^{n+1})$, by direct bounding).
\end{enumerate}
\end{method}

\begin{example}\label{ex:b1:taylor:tan}
Expansion of $\tan$ at order $5$. Write $\tan x =
\sin x \cdot \frac{1}{\cos x}$:
\[
\frac{1}{\cos x} = \frac{1}{1 - \bigl(\frac{x^2}{2} - \frac{x^4}{24}
+ o(x^5)\bigr)}
= 1 + \Bigl(\frac{x^2}{2} - \frac{x^4}{24}\Bigr)
+ \Bigl(\frac{x^2}{2}\Bigr)^{\!2} + o(x^5)
= 1 + \frac{x^2}{2} + \frac{5x^4}{24} + o(x^5),
\]
then
\[
\tan x = \Bigl(x - \frac{x^3}{6} + \frac{x^5}{120}\Bigr)
\Bigl(1 + \frac{x^2}{2} + \frac{5x^4}{24}\Bigr) + o(x^5)
= x + \frac{x^3}{3} + \frac{2x^5}{15} + o(x^5) .
\]
\end{example}

\section{Applications}

\begin{example}[Limits]\label{ex:b1:taylor:limits}
\[
\lim_{x \to 0} \frac{x - \sin x}{x^3}:
\qquad
x - \sin x = \frac{x^3}{6} + o(x^3) \sim \frac{x^3}{6},
\qquad\text{so the limit is } \frac16
\]
--- settling the question raised in \cref{exo:b1:functions:9}.
Likewise $\displaystyle\lim_{x\to0}\Bigl(\frac{\sin
x}{x}\Bigr)^{1/x^2}$: the logarithm is
\[
\frac{1}{x^2}\ln\Bigl(1 - \frac{x^2}{6} + o(x^2)\Bigr)
= \frac{1}{x^2}\Bigl(-\frac{x^2}{6} + o(x^2)\Bigr)
\longrightarrow -\frac16,
\qquad\text{limit } \eu^{-1/6}.
\]
\end{example}

\begin{proposition}[Local behavior]\label{prop:b1:taylor:local}
Suppose $f(x) = f(a) + c\,(x - a)^p + o\bigl((x-a)^p\bigr)$ with $c
\neq 0$ (first nonzero term after the constant; $p \geq 2$ at a
critical point).
\begin{itemize}
  \item If $p$ is even: $f$ has a local minimum at $a$ if $c > 0$, a
        local maximum if $c < 0$.
  \item If $p$ is odd: no extremum ($f - f(a)$ changes sign);
        if moreover the expansion starts after a linear term $f'(a)(x
        - a)$, the graph crosses its tangent: an \emph{inflection}.
\end{itemize}
\end{proposition}

\begin{proof}
Near $a$, $f(x) - f(a) = (x-a)^p\bigl(c + o(1)\bigr)$ has the sign
of $c\,(x-a)^p$: constant sign for $p$ even, changing for $p$ odd.
\end{proof}

\begin{example}[Asymptote by expansion]\label{ex:b1:taylor:asymptote}
As $x \to +\infty$,
\[
\sqrt{x^2 + x} = x\sqrt{1 + \tfrac1x}
= x\Bigl(1 + \frac{1}{2x} - \frac{1}{8x^2} + o\bigl(\tfrac{1}{x^2}\bigr)\Bigr)
= x + \frac12 - \frac{1}{8x} + o\bigl(\tfrac 1x\bigr):
\]
the line $y = x + \frac12$ is an asymptote, approached \emph{from
below} (the next term $-\frac{1}{8x}$ is negative).
\end{example}

\section{Exercises}

\begin{exercise}[$\star$]\label{exo:b1:taylor:1}
Give the expansions at $0$: $\eu^{2x}$ to order $3$;
$\;\ln(1 - x)$ to order $4$; $\;\sqrt{1 + x}$ to order $3$;
$\;\dfrac{1}{1 + x^2}$ to order $6$.
\end{exercise}

\begin{exercise}[$\star$]\label{exo:b1:taylor:2}
Compute the limits:
\[
\lim_{x\to 0} \frac{\eu^x - 1 - x}{x^2},
\qquad
\lim_{x\to 0} \frac{\cos x - \sqrt{1 - x^2}}{x^4},
\qquad
\lim_{x\to 0} \frac{\ln(1+x) - \sin x}{x^2}.
\]
\end{exercise}

\begin{exercise}[$\star$]\label{exo:b1:taylor:3}
Expand $\arctan x$ at $0$ to order $5$ by integrating the expansion
of $\frac{1}{1 + x^2}$, and $\arcsin x$ to order $5$ by integrating
that of $(1 - x^2)^{-1/2}$.
\end{exercise}

\begin{exercise}[$\star$]\label{exo:b1:taylor:4}
Using Taylor--Lagrange for $\exp$ on $\intcc{0}{1}$, prove that
\[
\Bigl| \eu - \sum_{k=0}^{n} \frac{1}{k!} \Bigr| \leq
\frac{3}{(n+1)!},
\]
and determine an $n$ guaranteeing $6$ exact decimals of $\eu$.
\end{exercise}

\begin{exercise}[$\star\star$]\label{exo:b1:taylor:5}
Expand to order $2$ in $\frac1n$ and deduce the limit and the
speed of convergence:
\[
\Bigl(1 + \frac 1n\Bigr)^{\!n}
= \eu\Bigl(1 - \frac{1}{2n} + \frac{11}{24n^2} +
o\Bigl(\frac{1}{n^2}\Bigr)\Bigr).
\]
\end{exercise}

\begin{exercise}[$\star\star$]\label{exo:b1:taylor:6}
Study the local behavior at $0$ of $f(x) = x^2 - x^4$ and of $g(x) =
x^3 + x^5$; and find the position of the graph of $h(x) = \eu^x$
relative to its tangent at $a = 1$, locally then globally.
\end{exercise}

\begin{exercise}[$\star\star$]\label{exo:b1:taylor:7}
Determine the asymptotes at $\pm\infty$ of $f(x) =
\sqrt[3]{x^3 + x^2}$ and the position of the curve relative to
them.
\end{exercise}

\begin{exercise}[$\star\star$]\label{exo:b1:taylor:8}
Find the equivalent, as $n \to \infty$, of
\[
u_n = \sqrt{n+1} - \sqrt n, \qquad
v_n = \ln(n+1) - \ln n, \qquad
w_n = \sin\frac{1}{n} - \tan\frac{1}{n},
\]
each as a power of $n$ times a constant.
\end{exercise}

\begin{exercise}[$\star\star\star$]\label{exo:b1:taylor:9}
Let $f$ be $C^2$ on $\R$. Prove that for every $x$ and $h > 0$:
\[
\abs{f'(x)} \leq \frac{\abs{f(x+h) - f(x-h)}}{2h} +
\frac{h}{2}\sup_{\intcc{x-h}{x+h}}\abs{f''} ,
\]
and deduce the Landau--Kolmogorov-type inequality: if $\abs f \leq
M_0$ and $\abs{f''} \leq M_2$ on $\R$, then $\abs{f'} \leq
\sqrt{2 M_0 M_2}$ everywhere. \emph{(Optimize over $h$.)}
\end{exercise}

\begin{exercise}[$\star\star\star$]\label{exo:b1:taylor:10}
The sequence $u_0 \in \intoo{0}{\pi}$, $u_{n+1} = \sin u_n$
decreases to $0$ (justify briefly). To find its speed, consider $v_n
= \frac{1}{u_n^2}$:
\begin{enumerate}
  \item using the expansion of $\sin$, prove $v_{n+1} - v_n \to
        \frac13$;
  \item with Cesàro (\cref{exo:b1:seq:10}), deduce $\frac{v_n}{n}
        \to \frac13$, then the equivalent $u_n \sim
        \sqrt{\dfrac{3}{n}}$.
\end{enumerate}
\end{exercise}
